% CS301 Week 9 -- Efficient sorting: divide-and-conquer and partition (merge
% sort, quick sort), the comparison-sort lower bound, and the three ways past
% the quadratic wall -- heap sort (guaranteed n log n, O(1) space), radix sort
% (leaves the comparison model), and shell sort (diminishing increments) --
% closing with a comparative analysis.
% Difficulty: intro (course-wide default per knowledge-base-CS301.md).
% Page-grounded in Sedgewick2011 Sec.2.2 (Mergesort, p.270), Sec.2.3
% (Quicksort, p.288), Sec.2.4 (Priority Queues / heapsort, p.308), Sec.5.1
% (String Sorts / radix, p.702), and the order-of-growth table p.187;
% Wirth2004 Sec.2.3.1 (Insertion Sort by Diminishing Increment, p.52),
% Sec.2.3.3 (Partition Sort, p.57), Sec.2.3.5 (Comparison of Array Sorting
% Methods, p.62), Sec.2.4.1 (Straight Merging, p.63); Karumanchi2017 Ch.10
% (Sorting, pp.504-546 PDF: merge/quick/shell/radix) and Ch.7 (Priority
% Queues and Heaps, pp.369-408 PDF, heapsort Sec.7.7). Aho-Hopcroft-Ullman
% Ch.8, Horowitz & Sahni Ch.7, Knuth Vol.3, and Neapolitan Ch.2 are cited at
% chapter level only (not page-indexed) -- no page numbers are asserted.
% Handoff: Week 8 closed on elementary sorting and promised "Why can no
% comparison sort beat Theta(n log n), and which algorithms escape that bound?
% Efficient sorting: quick, merge, shell, radix, heap." This deck pays that
% off. It reuses Week 2's O/Theta/Omega and recurrences, Week 3's array
% address arithmetic, Week 7's linked-list merge, and Week 8's stable/in-place
% vocabulary; it does not re-teach them. Continues the expression-processor
% thread: the calculator's 2,000-identifier table is sorted once at startup
% and searched many times, so the sorting cost is the startup pause to remove.
% Sets up Week 10 (file organization and hashing) and Week 11 (heaps are the
% array-as-tree idea that becomes the binary tree).

\documentclass[aspectratio=169]{beamer}
\input{header}
\coursecode{CS301}

\title{Beyond the Quadratic Wall}
\subtitle{Week 9 --- Efficient Sorting \& the Comparison Lower Bound}
\author[Dr. Auqib Hamid Lone]{\textbf{Dr. Auqib Hamid Lone}}
\institute[GCET Ganderbal]{PCC CS-301: Data Structures\\GCET, Ganderbal}
\date[Week 9]{\small Week 9}

\begin{document}

\frame{\titlepage}

% =============================================================================
% ACT 0 --- MOTIVATION + HANDOFF
% =============================================================================
\section{Motivation and the Handoff}

\begin{frame}{Today's Five Questions}
  \begin{enumerate}
    \item Merge and quick both promise $n \log n$ --- where does the $\log$
      come from, and when does quick break that promise?
    \item Is there a ceiling that \emph{no} comparison sort can pass?
    \item How does heap sort guarantee $n \log n$ with only $O(1)$ extra
      space?
    \item How does radix sort beat a bound that supposedly applies to every
      sort?
    \item Shell sort sits between insertion and quick --- when is it the
      right tool?
  \end{enumerate}
  \vspace{0.2cm}
  \begin{exampleblock}{Plan}
    Climb from quadratic to $n \log n$ twice --- once by dividing, once by
    partitioning. Prove the wall. Then go around it three ways: a guarantee,
    a changed model, and a better insertion sort.
  \end{exampleblock}
\end{frame}

\begin{frame}{Handoff: What Week 8 Left Open}
  \begin{itemize}
    \item Week 8 built the search contract (binary search is
      $\Theta(\log n)$ \emph{once the data is sorted}) and the four
      elementary sorts --- insertion, selection, bubble, counting --- all
      priced at $\Theta(n^2)$ worst case except counting.
    \item It closed on one promise: ``Why can no \emph{comparison} sort beat
      $\Theta(n \log n)$, and which algorithms escape that bound? Efficient
      sorting: quick, merge, shell, radix, heap.''
    \item The deficiency it left: the calculator's $2{,}000$-identifier
      table is sorted once at startup. Insertion sort pays about
      $2{,}000^2/2 = 2{,}000{,}000$ comparisons --- a visible pause before
      the first prompt.
    \item New symbols today: the recurrence $T(n) = 2T(n/2) + \Theta(n)$;
      the \key{pivot}; the \key{heap property}; the radix \key{digit} and the
      shell \key{gap}.
  \end{itemize}
  \begin{exampleblock}{This week's job}
    Turn that $2{,}000{,}000$ into roughly $2{,}000 \times 11 \approx
    22{,}000$ comparisons --- and learn exactly which assumption buys the
    difference.
  \end{exampleblock}
\end{frame}

\begin{frame}{The Bill for Order, at Two Sizes}
  \begin{center}
  {\footnotesize
  \begin{tabular}{rrrr}
    \toprule
    $n$ & $\dfrac{n(n-1)}{2}$ (insertion, worst) & $n\log_2 n$ & ratio \\
    \midrule
    $2{,}000$ & $1{,}999{,}000$ & $\approx 22{,}000$ & $\approx 91\times$ \\
    $1{,}000{,}000$ & $\approx 5\times 10^{11}$ & $\approx 2\times 10^{7}$ &
      $\approx 25{,}000\times$ \\
    \bottomrule
  \end{tabular}}
  \end{center}
  \begin{itemize}
    \item The elementary sorts are not \emph{wrong} --- they are priced for a
      different $n$. At $n = 2{,}000$ the pause is annoying; at
      $n = 10^6$ it is the difference between a program that runs and one
      that does not.
    \item Every efficient sort this week answers one question: \emph{what
      work can be avoided by looking at the whole input at once, instead of
      one adjacent pair at a time?}
  \end{itemize}
  \muted{Running application (Weeks 3--10): one calculator, rebuilt each
    week. Week 9 asks how its identifier table gets sorted fast enough to be
    worth sorting.}
\end{frame}

\transitionslide{Insertion sort is a fine sort. It is just a quadratic one.}

% =============================================================================
% ACT 1 --- DIVIDE AND CONQUER: MERGE SORT
% =============================================================================
\section{Merge Sort}

\sectiondivider{Section 1}{Divide and Conquer: Merge Sort}

\begin{frame}{Divide and Conquer: Three Questions}
  \begin{itemize}
    \item Sorting $2{,}000$ records in one go is hard. Sorting two
      $1{,}000$-record halves is easier. Combining two \emph{sorted} halves
      is easier still --- and costs one linear pass.
    \item That is the whole divide-and-conquer pattern: \key{divide} the
      input, \key{conquer} each part recursively, \key{combine} the answers.
  \end{itemize}
  \begin{block}{Definition --- divide-and-conquer sorting
    (Aho--Hopcroft--Ullman Ch.~8, standard treatment)}
    A sorting algorithm is divide-and-conquer when it (1) splits the input
    into independent subproblems, (2) sorts each recursively, and (3)
    combines the sorted parts without re-examining more than a constant
    number of times per element. The recurrence is then
    $T(n) = a\,T(n/b) + \Theta(\text{combine})$.
  \end{block}
  \begin{itemize}
    \item Merge sort answers the three questions: split at the midpoint;
      conquer two halves; combine by \emph{merging} two sorted runs. The
      combine step is the only new machinery.
  \end{itemize}
\end{frame}

\begin{frame}{The Merge Step: Two Sorted Runs, One Output}
  \begin{itemize}
    \item If both halves are already sorted, we do not need to sort again ---
      we need to \emph{interleave} them. Compare the two front elements, take
      the smaller, advance that run, repeat.
  \end{itemize}
  \begin{block}{Definition --- \texttt{merge(A, lo, mid, hi)}
    (Sedgewick Sec.~2.2, p.~270; Wirth Sec.~2.4.1, p.~63)}
    Precondition: $A[\texttt{lo}..\texttt{mid}]$ and
    $A[\texttt{mid}+1..\texttt{hi}]$ are each sorted. Postcondition:
    $A[\texttt{lo}..\texttt{hi}]$ is sorted. Ties take the \emph{left} run
    first --- that one comparison is what makes the merge stable.
  \end{block}
  \begin{itemize}
    \item Worked example: merge $[3,8,15]$ with $[2,9,11]$.
      Compare $3$ vs.\ $2$ $\to$ take $2$; $3$ vs.\ $9$ $\to$ take $3$;
      $8$ vs.\ $9$ $\to$ take $8$; $15$ vs.\ $9$ $\to$ take $9$;
      $15$ vs.\ $11$ $\to$ take $11$; drain $15$.
      Result: $[2,3,8,9,11,15]$.
    \item Cost: every element is copied once into the output, so the merge
      is $\Theta(n)$ on a range of $n$ elements.
  \end{itemize}
\end{frame}

\begin{frame}{Merge, Drawn}
  \begin{center}
    % Diagram: merging left run [3,8,15] with right run [2,9,11]. Each run is
    % a row of three cells; the output row below receives the smaller front
    % element. Coordinates: cells 0.9 cm wide at x = 0,1.1,2.2 (left) and
    % x = 3.6,4.7,5.8 (right); output cells at x = 0..5.5, y = -0.9.
    \begin{tikzpicture}[every node/.style={font=\scriptsize}]
      \foreach \i/\v in {0/3,1/8,2/15}
        \node[draw, fill=light-bg, minimum width=0.9cm,
          minimum height=0.6cm, align=center] (l\i) at (\i*1.1, 0.9) {\v};
      \foreach \i/\v in {0/2,1/9,2/11}
        \node[draw, fill=light-bg, minimum width=0.9cm,
          minimum height=0.6cm, align=center] (r\i) at (3.6+\i*1.1, 0.9) {\v};
      \node[anchor=east, text=neutral] at (-0.65, 0.9) {left run};
      \node[anchor=west, text=neutral] at (7.0, 0.9) {right run};
      \foreach \i/\v in {0/2,1/3,2/8,3/9,4/11,5/15}
        \node[draw, fill=primary-gold!20, minimum width=0.9cm,
          minimum height=0.6cm, align=center] (o\i) at (\i*1.1, -0.9) {\v};
      \node[anchor=east, text=neutral] at (-0.65, -0.9) {output};
      \draw[->, thick, primary-blue] (0.0, 0.45) -- (0.0, -0.35);
      \draw[->, thick, primary-blue] (3.6, 0.45) -- (0.0, -0.35);
      \node[text=primary-gold, align=center] at (2.9, -1.65)
        {compare the two fronts; take the smaller; repeat};
    \end{tikzpicture}
  \end{center}
  \vspace{-0.1cm}
  \begin{itemize}
    \item At most five comparisons, six output slots: the final element is
      drained without a comparison. The arrows show the first step --- front
      $3$ vs.\ front $2$, take $2$.
  \end{itemize}
\end{frame}

\begin{frame}[fragile]{The Merge Step in C}
  \footnotesize
  \begin{block}{Merging with one shared buffer (Sedgewick Sec.~2.2, p.~270;
    Wirth Sec.~2.4.1, p.~63)
    \cite{Sedgewick2011_algorithms,Wirth2004_algorithms_data_structures}}
\begin{semiverbatim}
void merge(int A[], int lo, int mid, int hi, int *tmp) {
  int i = lo, j = mid + 1, k = lo;
  while (i <= mid && j <= hi)
    tmp[k++] = (A[i] <= A[j]) ? A[i++] : A[j++];  // <= keeps it stable
  while (i <= mid) tmp[k++] = A[i++];             // drain the left run
  while (j <= hi)  tmp[k++] = A[j++];             // drain the right run
  for (k = lo; k <= hi; k++) A[k] = tmp[k];       // copy back
}
\end{semiverbatim}
  \end{block}
  \begin{itemize}
    \item The \texttt{<=} is load-bearing: on a tie, taking the \emph{left}
      element first preserves the original relative order --- stability, the
      Week 8 property, decided by one character.
    \item The buffer \texttt{tmp} is allocated once for the whole sort, not
      once per merge --- Week 1's discipline: one \texttt{malloc}, one
      \texttt{free}.
  \end{itemize}
\end{frame}

\begin{frame}[fragile]{Merge Sort in C}
  \scriptsize
  \begin{block}{Recursive merge sort (Sedgewick Sec.~2.2, p.~270; Karumanchi
    Ch.~10, pp.504--546 PDF)
    \cite{Sedgewick2011_algorithms,Karumanchi2017_data_structures_algorithms_made_easy}}
\begin{semiverbatim}
void msort(int A[], int lo, int hi, int *tmp) {
  if (lo >= hi) return;                 // 0 or 1 element is sorted
  int mid = lo + (hi - lo) / 2;
  msort(A, lo, mid, tmp);               // sort the left half
  msort(A, mid + 1, hi, tmp);           // sort the right half
  merge(A, lo, mid, hi, tmp);           // combine
}
void mergeSort(int A[], int n) {
  if (n < 2) return;
  int *tmp = malloc(n * sizeof(int));   // one buffer for the whole sort
  if (!tmp) return;                     // handle the allocation failure
  msort(A, 0, n - 1, tmp);
  free(tmp);
}
\end{semiverbatim}
  \end{block}
  \vspace{-0.25cm}
  \begin{itemize}
    \item Recursion depth $\log_2 n$: the \emph{stack} costs $\Theta(\log n)$;
      the \emph{heap} buffer $\Theta(n)$ --- two memories, two bills.
  \end{itemize}
\end{frame}

\begin{frame}{The Recurrence and Its Tree}
  \begin{itemize}
    \item Let $T(n)$ be the worst-case cost of sorting $n$ elements. Split
      costs $\Theta(1)$; two halves of size $n/2$; merge costs $\Theta(n)$:
      \[ T(1) = \Theta(1), \qquad T(n) = 2\,T(n/2) + \Theta(n). \]
    \item Unroll it as a tree. Every level does $\Theta(n)$ total work ---
      level $0$ merges one range of $n$, level $1$ merges two ranges of
      $n/2$, and so on. There are $\log_2 n + 1$ levels, so
      $T(n) = \Theta(n \log n)$.
  \end{itemize}
  \begin{center}
    % Diagram: recursion tree for T(n) = 2T(n/2) + cn. Root at (0,3.2);
    % children at (+-2.2,2.0); grandchildren at (+-3.2,+-1.2,0.8). Edges are
    % drawn between named nodes so they stop at the circle borders instead of
    % cutting through the labels. The right margin annotates the cn per level.
    \begin{tikzpicture}[every node/.style={font=\scriptsize}]
      \node[circle, draw=primary-blue, fill=light-bg, inner sep=2pt]
        (t0) at (0, 3.2) {$T(n)$};
      \node[circle, draw=primary-blue, fill=light-bg, inner sep=2pt]
        (t1) at (-2.2, 2.0) {$T(n/2)$};
      \node[circle, draw=primary-blue, fill=light-bg, inner sep=2pt]
        (t2) at (2.2, 2.0) {$T(n/2)$};
      \foreach \x/\lab/\nm in {-3.2/{$T(n/4)$}/g1,-1.2/{$T(n/4)$}/g2,
                                1.2/{$T(n/4)$}/g3,3.2/{$T(n/4)$}/g4}
        \node[circle, draw=primary-blue, fill=light-bg, inner sep=1.5pt]
          (\nm) at (\x, 0.8) {\lab};
      \draw[thick, primary-blue] (t0) -- (t1);
      \draw[thick, primary-blue] (t0) -- (t2);
      \draw[thick, primary-blue] (t1) -- (g1);
      \draw[thick, primary-blue] (t1) -- (g2);
      \draw[thick, primary-blue] (t2) -- (g3);
      \draw[thick, primary-blue] (t2) -- (g4);
      \node[text=neutral, anchor=west] at (3.9, 3.2) {level $0$: $cn$};
      \node[text=neutral, anchor=west] at (3.9, 2.0) {level $1$: $2\cdot c(n/2) = cn$};
      \node[text=neutral, anchor=west] at (3.9, 0.8) {level $2$: $cn$};
      \node[text=primary-gold, align=center] at (0, -0.4)
        {$\log_2 n + 1$ levels, each $\Theta(n)$ $\;\Rightarrow\;$
          $T(n) = \Theta(n\log n)$};
    \end{tikzpicture}
  \end{center}
  \vspace{-0.35cm}
\end{frame}

\begin{frame}{Merge Sort, Drawn: Split Down, Merge Up}
  \begin{center}
    % Diagram: merge sort of [38,27,43,3,9,82,10,19]. Four rows show the
    % split from one 8-cell array to eight singleton arrays; the merge phase
    % runs the same ladder upward. Cells 0.8 cm wide at x = 0..7.
    \begin{tikzpicture}[every node/.style={font=\scriptsize}]
      % row 1: whole array
      \draw[rounded corners, fill=light-bg, draw=neutral!40]
        (-0.45, 2.25) rectangle (7.45, 2.85);
      \foreach \i/\v in {0/38,1/27,2/43,3/3,4/9,5/82,6/10,7/19}
        \node at (\i, 2.55) {\v};
      % row 2: two halves of four
      \draw[rounded corners, fill=light-bg, draw=neutral!40]
        (-0.45, 1.45) rectangle (3.45, 2.05);
      \draw[rounded corners, fill=light-bg, draw=neutral!40]
        (3.55, 1.45) rectangle (7.45, 2.05);
      \foreach \i/\v in {0/38,1/27,2/43,3/3,4/9,5/82,6/10,7/19}
        \node at (\i, 1.75) {\v};
      % row 3: four pairs
      \foreach \s in {0,2,4,6}
        \draw[rounded corners, fill=light-bg, draw=neutral!40]
          (\s-0.45, 0.65) rectangle (\s+1.45, 1.25);
      \foreach \i/\v in {0/38,1/27,2/43,3/3,4/9,5/82,6/10,7/19}
        \node at (\i, 0.95) {\v};
      % row 4: singletons
      \foreach \i/\v in {0/38,1/27,2/43,3/3,4/9,5/82,6/10,7/19}
        \node[draw, fill=white, minimum width=0.72cm, minimum height=0.5cm,
          inner sep=0] at (\i, 0.05) {\v};
      % merge-up arrow on the right
      \draw[->, thick, primary-gold] (7.9, -0.15) -- (7.9, 2.85);
      \node[text=primary-gold, rotate=90, anchor=south] at (8.45, 1.35)
        {merge up};
      \node[text=neutral, anchor=east] at (-0.7, 2.85) {\tiny split};
    \end{tikzpicture}
  \end{center}
  \vspace{-0.1cm}
  \begin{itemize}
    \item Splitting is free (index arithmetic); all the work is in the
      upward merges. Each level merges every element once, which is the
      $cn$ per level in the recurrence.
  \end{itemize}
\end{frame}

\begin{frame}{Merge Sort: The Bill}
  \begin{center}
  {\footnotesize
  \begin{tabular}{lll}
    \toprule
    \textbf{case} & \textbf{comparisons} & \textbf{note} \\
    \midrule
    best & $\Theta(n\log n)$ & even sorted input must be merged \\
    average & $\Theta(n\log n)$ & no input shape helps or hurts \\
    worst & $\Theta(n\log n)$ & predictable --- no adversarial input \\
    \bottomrule
  \end{tabular}}
  \end{center}
  \begin{itemize}
    \item \good{Stable} (ties take the left run) and \good{predictable}: the
      cost does not depend on the input's order. That predictability is worth
      real money in a library sort.
    \item \bad{Not in-place}: a $\Theta(n)$ auxiliary buffer, plus
      $\Theta(\log n)$ recursion stack.
    \item \good{The right tool for linked lists}: merging two sorted lists
      needs no buffer at all --- just pointer splicing, which is why Week 7's
      list merge is the model for external sorting in Week 10.
  \end{itemize}
  \muted{Sedgewick's standard implementation uses about $n\log_2 n$ compares
    in the worst case (Sec.~2.2, p.~270); the exact count differs by a small
    constant, the $\Theta$ class does not.}
\end{frame}

\transitionslide{Merge sort guarantees the bound. Can we reach it without the buffer?}

% =============================================================================
% ACT 2 --- PARTITION: QUICK SORT
% =============================================================================
\section{Quick Sort}

\sectiondivider{Section 2}{Partition: Quick Sort}

\begin{frame}{The Partition Idea}
  \begin{itemize}
    \item Merge sort does its work \emph{after} the recursion, in the merge.
      Quick sort does its work \emph{before}: choose one element as a
      \key{pivot}, rearrange the array so everything smaller sits left of it
      and everything larger sits right, then recurse on the two sides.
    \item There is no merge step. Once the pivot is placed, it is in its
      \emph{final} position and never moves again.
  \end{itemize}
  \begin{block}{Definition --- partition (Wirth Sec.~2.3.3, p.~57;
    Sedgewick Sec.~2.3, p.~288)}
    Choose a pivot $p$. Rearrange $A[\texttt{lo}..\texttt{hi}]$ so that some
    index $q$ holds $p$, every element in $A[\texttt{lo}..q-1]$ is $\le p$,
    and every element in $A[q+1..\texttt{hi}]$ is $\ge p$. Return $q$. The
    partition costs one linear pass, $\Theta(n)$, with $O(1)$ extra space.
  \end{block}
  \begin{itemize}
    \item Then: quick sort left of $q$, quick sort right of $q$. The
      recurrence has the same \emph{shape} as merge sort's,
      $T(n) = T(k) + T(n-1-k) + \Theta(n)$, but the split $k$ is now
      data-dependent --- and that is where the danger lives.
  \end{itemize}
\end{frame}

\begin{frame}[fragile]{Lomuto Partition in C}
  \footnotesize
  \begin{block}{Partition around the last element (Wirth Sec.~2.3.3, p.~57;
    Sedgewick Sec.~2.3, p.~288)
    \cite{Wirth2004_algorithms_data_structures,Sedgewick2011_algorithms}}
\begin{semiverbatim}
int partition(int A[], int lo, int hi) {
  int pivot = A[hi];            // last element as the pivot
  int i = lo - 1;               // boundary of the <= pivot region
  for (int j = lo; j < hi; j++)
    if (A[j] <= pivot) {        // <= keeps equal keys on the left
      i++;
      swap(&A[i], &A[j]);
    }
  swap(&A[i + 1], &A[hi]);      // pivot into its final slot
  return i + 1;
}
\end{semiverbatim}
  \end{block}
  \begin{itemize}
    \item Invariant before each iteration: $A[\texttt{lo}..i] \le$ pivot,
      $A[i+1..j-1] >$ pivot, and $A[\texttt{hi}]$ holds the pivot.
      One pass, $O(1)$ extra space.
  \end{itemize}
\end{frame}

\begin{frame}{Partition, Drawn}
  \begin{center}
    % Diagram: Lomuto partition of [7,2,9,4,3,8,5] with pivot 5.
    % Three rows: initial, after j=4, final. Cells at x = 0..6 (1 cm apart),
    % rows y = 1.5, 0.4, -0.7. Green = <= pivot region, gold = pivot.
    \begin{tikzpicture}[every node/.style={font=\scriptsize}]
      \foreach \i/\v in {0/7,1/2,2/9,3/4,4/3,5/8,6/5}
        \node[draw, fill=white, minimum width=0.8cm, minimum height=0.6cm,
          align=center] at (\i, 1.5) {\v};
      \node[draw, fill=primary-gold!30, minimum width=0.8cm,
        minimum height=0.6cm, align=center] at (6, 1.5) {5};
      \node[anchor=east, text=neutral] at (-0.6, 1.5) {start};
      \node[text=negative, anchor=south] at (6, 1.85) {\tiny pivot};
      \foreach \i/\v in {0/2,1/4,2/3,3/7,4/9,5/8,6/5}
        \node[draw, fill=white, minimum width=0.8cm, minimum height=0.6cm,
          align=center] at (\i, 0.4) {\v};
      \foreach \i in {0,1,2}
        \node[draw, fill=positive!25, minimum width=0.8cm,
          minimum height=0.6cm, align=center] at (\i, 0.4) {\tiny $\le 5$};
      \node[draw, fill=primary-gold!30, minimum width=0.8cm,
        minimum height=0.6cm, align=center] at (6, 0.4) {5};
      \node[anchor=east, text=neutral] at (-0.6, 0.4) {after $j=4$};
      \foreach \i/\v in {0/2,1/4,2/3,3/5,4/9,5/8,6/7}
        \node[draw, fill=white, minimum width=0.8cm, minimum height=0.6cm,
          align=center] at (\i, -0.7) {\v};
      \foreach \i in {0,1,2}
        \node[draw, fill=positive!25, minimum width=0.8cm,
          minimum height=0.6cm, align=center] at (\i, -0.7) {\tiny $\le 5$};
      \node[draw, fill=primary-gold!30, minimum width=0.8cm,
        minimum height=0.6cm, align=center] at (3, -0.7) {5};
      \node[anchor=east, text=neutral] at (-0.6, -0.7) {final};
      \node[text=primary-gold, align=center] at (3.0, -1.45)
        {pivot at index $3$: left $\le 5$, right $> 5$ --- final position};
    \end{tikzpicture}
  \end{center}
  \vspace{-0.1cm}
  \begin{itemize}
    \item Scanning $7,2,9,4,3,8$ moves each element $\le 5$ into the green
      region. The last swap puts $5$ between the regions, at index $3$.
  \end{itemize}
\end{frame}

\begin{frame}[fragile]{Quick Sort in C}
  \footnotesize
  \begin{block}{Quick sort with a bounded recursion stack (Sedgewick
    Sec.~2.3, p.~288; Karumanchi Ch.~10, pp.504--546 PDF)
    \cite{Sedgewick2011_algorithms,Karumanchi2017_data_structures_algorithms_made_easy}}
\begin{semiverbatim}
void qsortRange(int A[], int lo, int hi) {
  while (lo < hi) {                 // loop = tail-call elimination
    int p = partition(A, lo, hi);
    if (p - lo < hi - p) {          // recurse on the SMALLER side
      qsortRange(A, lo, p - 1);
      lo = p + 1;                   // iterate on the larger side
    } else {
      qsortRange(A, p + 1, hi);
      hi = p - 1;
    }
  }
}
\end{semiverbatim}
  \end{block}
  \begin{itemize}
    \item Recursing on the smaller side first bounds the stack depth to
      $O(\log n)$ even when the splits are badly unbalanced.
  \end{itemize}
\end{frame}

\begin{frame}{Best and Average: Balanced Partitions}
  \begin{itemize}
    \item If every pivot lands near the middle, the two sides are about
      $n/2$ each and the recurrence is the merge-sort recurrence:
      \[ T(n) = 2\,T(n/2) + \Theta(n) = \Theta(n\log n). \]
    \item The average case is also $\Theta(n\log n)$ over random pivots.
      More strongly, the splits need not be even: with a $9{:}1$ split at
      every level,
      \[ T(n) = T(9n/10) + T(n/10) + \Theta(n) = \Theta(n\log n). \]
    \item The reason: each level still does $\Theta(n)$ total work; only the
      \emph{number of levels} grows, and it stays $\Theta(\log n)$ as long as
      every split removes a constant fraction. Constant-factor imbalance is
      harmless --- a split that removes only \emph{one} element is not.
  \end{itemize}
  \begin{exampleblock}{The takeaway}
    Quick sort's $\log$ comes from halving the \emph{problem count}, not from
    the pivot being the median. Any split that keeps shrinking the pieces
    geometrically keeps the $\log$.
  \end{exampleblock}
\end{frame}

\begin{frame}{Worst Case: Sorted Input, Last-Element Pivot}
  \begin{itemize}
    \item Take the sorted array $[1,2,3,4,5,6]$ and the code above. The
      pivot is always \texttt{A[hi]}, which is always the \emph{largest}
      remaining element.
    \item The partition then produces a left side of $n-1$ elements and a
      right side of $0$. The recurrence collapses to
      \[ T(n) = T(n-1) + \Theta(n) = \Theta(n^2), \]
      because the levels are a chain of lengths $n, n-1, \dots, 1$, not a
      tree of depth $\log n$.
  \end{itemize}
  \begin{center}
    % Diagram: the chain of partitions for sorted input. Each row shows the
    % live range; the gold cell is the pivot, placed at the right end. Cells
    % at x = 0..5, rows y = 1.4, 0.7, 0.0, -0.7; only the first three rows
    % are labelled.
    \begin{tikzpicture}[every node/.style={font=\scriptsize}]
      % row n=6: live cells 0..5, pivot at 5
      \draw[rounded corners, fill=light-bg, draw=neutral!40]
        (-0.38, 1.18) rectangle (4.62, 1.62);
      \draw[rounded corners, fill=primary-gold!30, draw=neutral!60]
        (4.65, 1.18) rectangle (5.38, 1.62);
      \node[anchor=east, text=neutral] at (-0.55, 1.4) {$n=6$};
      % row n=5: live cells 0..4, pivot at 4
      \draw[rounded corners, fill=light-bg, draw=neutral!40]
        (-0.38, 0.48) rectangle (3.62, 0.92);
      \draw[rounded corners, fill=primary-gold!30, draw=neutral!60]
        (3.65, 0.48) rectangle (4.38, 0.92);
      \node[anchor=east, text=neutral] at (-0.55, 0.7) {$n=5$};
      % row n=4: live cells 0..3, pivot at 3
      \draw[rounded corners, fill=light-bg, draw=neutral!40]
        (-0.38, -0.22) rectangle (2.62, 0.22);
      \draw[rounded corners, fill=primary-gold!30, draw=neutral!60]
        (2.65, -0.22) rectangle (3.38, 0.22);
      \node[anchor=east, text=neutral] at (-0.55, 0.0) {$n=4$};
      % row n=3: live cells 0..2, pivot at 2
      \draw[rounded corners, fill=light-bg, draw=neutral!40]
        (-0.38, -0.92) rectangle (1.62, -0.48);
      \draw[rounded corners, fill=primary-gold!30, draw=neutral!60]
        (1.65, -0.92) rectangle (2.38, -0.48);
      \node[anchor=east, text=neutral] at (-0.55, -0.7) {$n=3$};
      \node[text=primary-gold, align=center] at (2.5, -1.55)
        {each partition removes one pivot: $n + (n-1) + \cdots + 1$ work};
    \end{tikzpicture}
  \end{center}
\end{frame}

\begin{frame}{Choosing a Pivot}
  \begin{itemize}
    \item The worst case is not exotic: \emph{sorted input} is exactly what
      our identifier table arrives as after the tokenizer has run. A naive
      quicksort would undo the whole point of sorting.
    \item Two standard defences:
      \begin{enumerate}
        \item \key{Median-of-three}: take the median of $A[\texttt{lo}]$,
          $A[\text{mid}]$, $A[\texttt{hi}]$ as pivot. On sorted or
          reverse-sorted input this picks the middle and restores balanced
          splits --- the two most common adversarial shapes are removed.
        \item \key{Randomized pivot}: swap a uniformly random index into
          \texttt{hi} before partitioning. No fixed input can force the worst
          case; the expected cost over the algorithm's own randomness is
          $\Theta(n\log n)$.
      \end{enumerate}
    \item Neither removes the $\Theta(n^2)$ worst case from the
      \emph{adversary's} model --- but randomization makes hitting it
      astronomically unlikely, and median-of-three makes the natural bad
      inputs good.
  \end{itemize}
  \muted{Practical libraries go one step further: \emph{introsort} runs quick
    sort, tracks recursion depth, and switches to heap sort if the depth
    exceeds $2\log_2 n$ --- quick sort's speed with heap sort's guarantee.}
\end{frame}

\begin{frame}{Quick Sort: The Bill}
  \begin{center}
  {\footnotesize
  \begin{tabular}{lll}
    \toprule
    \textbf{case} & \textbf{comparisons} & \textbf{example input} \\
    \midrule
    best & $\Theta(n\log n)$ & balanced pivots \\
    average & $\Theta(n\log n)$ & random order \\
    worst & $\Theta(n^2)$ & sorted / reverse-sorted, last-element pivot \\
    \bottomrule
  \end{tabular}}
  \end{center}
  \begin{itemize}
    \item \good{In-place}: $O(1)$ extra space for the partition, $O(\log n)$
      recursion stack with the smaller-side rule. No $\Theta(n)$ buffer.
    \item \bad{Not stable}: the long-range swaps leap equal keys over each
      other, exactly as selection sort did in Week 8.
    \item \good{Cache-friendly}: the partition scans the array sequentially,
      which is why quick sort usually beats merge sort in wall-clock time
      despite the same $\Theta$ class --- and why it is the default
      general-purpose sort in most libraries.
  \end{itemize}
\end{frame}

\begin{frame}{Check Your Prediction: Sorted Input}
  \begin{block}{You quick sort $[1,2,3,4,5,6]$ with the last-element pivot.
    How many comparisons does \texttt{partition} perform in total?}
    A: $15$ \quad B: $6$ \quad C: $30$
  \end{block}
  \begin{enumerate}
    \item A --- the first partition scans $5$ elements, the next $4$, then
      $3, 2, 1$: $5+4+3+2+1 = 15 = n(n-1)/2$.
    \item B --- that is only one comparison per element, the cost of a single
      pass; it ignores the $n-1$ recursive partitions.
    \item C --- that is $n(n-1)$, double-counting each partition.
  \end{enumerate}
  \muted{Correct: A. The sorted input forces the worst case; the total is
    $\Theta(n^2)$. With median-of-three on the same input the pivot is the
    middle and the count drops to $\Theta(n\log n)$.}
\end{frame}

\transitionslide{Two sorts reach $n\log n$. Is that a coincidence, or a ceiling?}

% =============================================================================
% ACT 3 --- THE COMPARISON LOWER BOUND
% =============================================================================
\section{The Comparison Lower Bound}

\sectiondivider{Section 3}{The Comparison-Sort Lower Bound}

\begin{frame}{Count Comparisons, Not Seconds}
  \begin{itemize}
    \item Wall-clock time depends on the machine. The portable question is:
      how many \emph{comparisons} must a sort make before it can be sure of
      the order?
    \item Every comparison sort --- insertion, merge, quick, heap --- learns
      about the data only by asking questions of the form ``is $A[i] <
      A[j]$?''. Each answer is yes or no.
    \item So a run on a fixed input follows one path from the root of a
      \key{decision tree} to a leaf, where each internal node is one
      comparison and each leaf is the output permutation the algorithm
      believes.
  \end{itemize}
  \begin{block}{The model (Aho--Hopcroft--Ullman Ch.~8, standard treatment;
    Knuth Vol.~3, standard treatment)}
    Only key comparisons are counted. Moving, copying, or indexing is free.
    A correct comparison sort must reach a \emph{different} leaf for each of
    the $n!$ possible input orderings --- otherwise two orderings that end at
    the same leaf get the same output, and one of them is wrong.
    \cite{AhoHopcroftUllman1983_data_structures_algorithms,Knuth1998_taocp_volume3}
  \end{block}
\end{frame}

\begin{frame}{The Decision Tree}
  \begin{center}
    % A full binary decision tree with six leaves, the count a correct
    % comparison sort of n=3 must distinguish. Every internal node has two
    % children (a comparison has two outcomes). Root at (0,3.0); height 3.
    \begin{tikzpicture}[every node/.style={font=\scriptsize}]
      \node[circle, draw=primary-blue, fill=light-bg, inner sep=1.5pt]
        (r) at (0, 3.0) {};
      \node[circle, draw=primary-blue, fill=light-bg, inner sep=1.5pt]
        (l) at (-2.0, 2.1) {};
      \node[circle, draw=primary-blue, fill=light-bg, inner sep=1.5pt]
        (rr) at (2.0, 2.1) {};
      \node[draw, rounded corners, fill=positive!15, inner sep=2pt]
        (l1) at (-3.0, 1.2) {$\pi_1$};
      \node[circle, draw=primary-blue, fill=light-bg, inner sep=1.5pt]
        (l2) at (-1.0, 1.2) {};
      \node[draw, rounded corners, fill=positive!15, inner sep=2pt]
        (r1) at (1.0, 1.2) {$\pi_4$};
      \node[circle, draw=primary-blue, fill=light-bg, inner sep=1.5pt]
        (r2) at (3.0, 1.2) {};
      \node[draw, rounded corners, fill=positive!15, inner sep=2pt]
        (l21) at (-1.5, 0.3) {$\pi_2$};
      \node[draw, rounded corners, fill=positive!15, inner sep=2pt]
        (l22) at (-0.5, 0.3) {$\pi_3$};
      \node[draw, rounded corners, fill=positive!15, inner sep=2pt]
        (r21) at (2.5, 0.3) {$\pi_5$};
      \node[draw, rounded corners, fill=positive!15, inner sep=2pt]
        (r22) at (3.5, 0.3) {$\pi_6$};
      \draw[thick, primary-blue] (r) -- node[fill=white, inner sep=1pt,
        font=\tiny, text=primary-blue] {$\le$} (l);
      \draw[thick, primary-blue] (r) -- node[fill=white, inner sep=1pt,
        font=\tiny, text=primary-blue] {$>$} (rr);
      \draw[thick, primary-blue] (l) -- (l1);
      \draw[thick, primary-blue] (l) -- (l2);
      \draw[thick, primary-blue] (rr) -- (r1);
      \draw[thick, primary-blue] (rr) -- (r2);
      \draw[thick, primary-blue] (l2) -- (l21);
      \draw[thick, primary-blue] (l2) -- (l22);
      \draw[thick, primary-blue] (r2) -- (r21);
      \draw[thick, primary-blue] (r2) -- (r22);
      \node[text=primary-gold, align=center] at (0.2, -0.7)
        {$n=3$: $3! = 6$ leaves, so height $\ge \lceil\log_2 6\rceil = 3$};
    \end{tikzpicture}
  \end{center}
  \vspace{-0.1cm}
  \begin{itemize}
    \item A binary tree with $L$ leaves has height at least
      $\lceil\log_2 L\rceil$. With $L = n!$ leaves, the worst-case number of
      comparisons is at least $\log_2(n!)$.
  \end{itemize}
\end{frame}

\begin{frame}{The Bound: $\log_2(n!) = \Theta(n\log n)$}
  \begin{itemize}
    \item A lower bound on $\log_2(n!)$ is enough. Keep only the larger half
      of the factors:
      \[ \log_2(n!) = \sum_{k=1}^{n}\log_2 k
         \;\ge\; \sum_{k=\lceil n/2\rceil}^{n}\log_2\!\left(\frac{n}{2}\right)
         \;=\; \frac{n}{2}\log_2\!\left(\frac{n}{2}\right)
         \;=\; \Omega(n\log n). \]
    \item The matching upper bound is immediate, $\log_2(n!) \le n\log_2 n$,
      so $\log_2(n!) = \Theta(n\log n)$.
    \item Stirling's approximation says the same thing more sharply:
      \[ \log_2(n!) = n\log_2 n - \frac{n}{\ln 2} + O(\log n)
         \approx n\log_2 n - 1.44\,n. \]
  \end{itemize}
  \begin{alertblock}{The comparison-sort lower bound}
    Every comparison sort makes at least $\Omega(n\log n)$ comparisons in the
    worst case. Merge sort and heap sort \emph{reach} this bound, which is
    why they are optimal among comparison sorts.
  \end{alertblock}
\end{frame}

\begin{frame}{What the Wall Does and Does Not Say}
  \begin{itemize}
    \item It is a \key{worst-case} bound over all inputs. An adaptive sort
      can be $O(n)$ on a particular input --- insertion sort on sorted data
      does $n-1$ comparisons --- but \emph{some} input still forces
      $\Omega(n\log n)$.
    \item It applies only to algorithms that learn about keys \emph{by
      comparing them}. It says nothing about a sort that reads the bits of a
      key directly.
    \item It does not say $n\log n$ is the fastest possible sorting time ---
      only the fastest possible \emph{comparison} sorting time.
  \end{itemize}
  \begin{center}
  {\footnotesize
  \begin{tabular}{rrrr}
    \toprule
    $n$ & $n!$ & $\log_2(n!)$ & $n\log_2 n$ \\
    \midrule
    $10$ & $3{,}628{,}800$ & $\approx 22$ & $\approx 33$ \\
    $20$ & $\approx 2.4\times10^{18}$ & $\approx 61$ & $\approx 86$ \\
    \bottomrule
  \end{tabular}}
  \end{center}
  \muted{The table is a sanity check, not a proof: the bound is asymptotic,
    and the constants above are not tight for small $n$.}
\end{frame}

\transitionslide{To go faster than the wall, stop comparing whole keys.}

% =============================================================================
% ACT 4 --- HEAP, RADIX, AND SHELL SORT
% =============================================================================
\section{Heap, Radix, and Shell Sort}

\sectiondivider{Section 4}{Three Ways Past the Wall}

\begin{frame}{Heap Sort: An Array That Thinks It Is a Tree}
  \begin{itemize}
    \item Merge sort has the guarantee but pays $\Theta(n)$ space. Quick sort
      has the space but loses the guarantee. Can we have both?
    \item Yes --- if we store a \emph{complete binary tree} in the array
      itself, with no pointers. A \key{max-heap} is such a tree where every
      parent is at least as large as its children.
  \end{itemize}
  \begin{block}{Definition --- binary max-heap (Karumanchi Ch.~7,
    pp.369--408 PDF, Sec.~7.7; Sedgewick Sec.~2.4, p.~308)}
    A binary max-heap is a complete binary tree with the \key{heap property}:
    for every node $i$, $A[i] \ge A[\text{left}(i)]$ and
    $A[i] \ge A[\text{right}(i)]$. In a 0-based array,
    $\text{left}(i)=2i+1$, $\text{right}(i)=2i+2$,
    $\text{parent}(i)=\lfloor(i-1)/2\rfloor$. The maximum is always
    $A[0]$.
  \end{block}
  \begin{itemize}
    \item No pointers, no allocation: the tree structure is pure index
      arithmetic, the same address arithmetic as Week 3's arrays. That is
      what makes the $O(1)$ space possible.
  \end{itemize}
\end{frame}

\begin{frame}{Heap Sort, Drawn: The Same Seven Values, Two Views}
  \begin{center}
    % Diagram: max-heap [50,30,40,10,20,35,25] drawn as a tree over the same
    % array. Tree root (0,2.5), children (+-1.6,1.6), grandchildren
    % (-2.4,-0.8,0.8,2.4 at y=0.7). Array cells at x = 0..6, y = -0.6.
    \begin{tikzpicture}[every node/.style={font=\scriptsize}]
      \node[circle, draw=primary-blue, fill=light-bg, inner sep=2pt]
        (n0) at (0, 2.5) {50};
      \node[circle, draw=primary-blue, fill=light-bg, inner sep=2pt]
        (n1) at (-1.6, 1.6) {30};
      \node[circle, draw=primary-blue, fill=light-bg, inner sep=2pt]
        (n2) at (1.6, 1.6) {40};
      \node[circle, draw=primary-gold, fill=light-bg, inner sep=2pt]
        (n3) at (-2.4, 0.7) {10};
      \node[circle, draw=primary-gold, fill=light-bg, inner sep=2pt]
        (n4) at (-0.8, 0.7) {20};
      \node[circle, draw=primary-gold, fill=light-bg, inner sep=2pt]
        (n5) at (0.8, 0.7) {35};
      \node[circle, draw=primary-gold, fill=light-bg, inner sep=2pt]
        (n6) at (2.4, 0.7) {25};
      \foreach \a/\b in {n0/n1,n0/n2,n1/n3,n1/n4,n2/n5,n2/n6}
        \draw[thick, primary-blue] (\a) -- (\b);
      \foreach \i/\v in {0/50,1/30,2/40,3/10,4/20,5/35,6/25}
        \node[draw, fill=light-bg, minimum width=0.72cm,
          minimum height=0.55cm, inner sep=0] at (\i, -0.6) {\v};
      \foreach \i in {0,...,6}
        \node[text=neutral] at (\i, -1.15) {\tiny \i};
      \node[text=neutral, anchor=east] at (-0.55, -0.6) {array};
      \node[text=primary-gold, align=center] at (4.05, 1.6)
        {gold leaves: no children,\\ so no sift-down work};
    \end{tikzpicture}
  \end{center}
  \vspace{-0.1cm}
  \begin{itemize}
    \item Check the property: $50 \ge 30,40$; $30 \ge 10,20$;
      $40 \ge 35,25$. The maximum sits at index $0$, which is why extracting
      it is cheap to locate --- the work is in restoring the property.
  \end{itemize}
\end{frame}

\begin{frame}[fragile]{Sift Down: Restore the Heap Property}
  \footnotesize
  \begin{block}{Sift-down (heapify) in $O(\log n)$ (Karumanchi Ch.~7,
    pp.369--408 PDF; Sedgewick Sec.~2.4, p.~308)
    \cite{Karumanchi2017_data_structures_algorithms_made_easy,Sedgewick2011_algorithms}}
\begin{semiverbatim}
void siftDown(int A[], int i, int n) {
  for (;;) {
    int l = 2 * i + 1, r = l + 1, m = i;
    if (l < n && A[l] > A[m]) m = l;   // larger child
    if (r < n && A[r] > A[m]) m = r;
    if (m == i) return;                // heap property restored
    swap(&A[i], &A[m]);
    i = m;                             // keep sinking
  }
}
\end{semiverbatim}
  \end{block}
  \begin{itemize}
    \item Only one subtree can violate the property after a swap, so each
      step follows a single root-to-leaf path: at most the tree height,
      $\lfloor\log_2 n\rfloor$ swaps.
  \end{itemize}
\end{frame}

\begin{frame}{Build a Heap in $\Theta(n)$, Not $\Theta(n\log n)$}
  \begin{itemize}
    \item The obvious build is $n$ insertions, each sifting up $O(\log n)$:
      $O(n\log n)$. It works, but it is not tight.
    \item The tight build sifts \emph{down} every internal node, from the last
      parent $\lfloor n/2\rfloor - 1$ back to index $0$. Leaves need no work
      at all.
    \item Count the work by height. About $n/2$ nodes are leaves (height
      $0$), $n/4$ have height $1$, $n/8$ height $2$, and so on:
      \[ \sum_{h \ge 0} \frac{n}{2^{h+1}}\,h
         \;=\; \frac{n}{2}\sum_{h\ge 0}\frac{h}{2^{h}}
         \;=\; \frac{n}{2}\cdot 2 \;=\; \Theta(n). \]
  \end{itemize}
  \begin{exampleblock}{Why most nodes are cheap}
    Half the nodes never move, a quarter move at most one level, an eighth at
    most two. The expensive top-of-tree nodes are few. This is the classic
    case where the obvious bound is off by a log factor --- the same
    amortized-counting habit Week 2 built.
  \end{exampleblock}
\end{frame}

\begin{frame}[fragile]{Heap Sort in C}
  \footnotesize
  \begin{block}{Heap sort (Karumanchi Ch.~7, pp.369--408 PDF, Sec.~7.7;
    Sedgewick Sec.~2.4, p.~308)
    \cite{Karumanchi2017_data_structures_algorithms_made_easy,Sedgewick2011_algorithms}}
\begin{semiverbatim}
void heapSort(int A[], int n) {
  for (int i = n / 2 - 1; i >= 0; i--)   // build max-heap: Theta(n)
    siftDown(A, i, n);
  for (int end = n - 1; end > 0; end--) {
    swap(&A[0], &A[end]);                // max to its final slot
    siftDown(A, 0, end);                 // restore on the shrinking heap
  }
}
\end{semiverbatim}
  \end{block}
  \begin{itemize}
    \item Phase 1 is $\Theta(n)$ by the height sum. Phase 2 performs $n-1$
      extractions, each $O(\log n)$: $\Theta(n\log n)$.
    \item Total: $\Theta(n\log n)$ in the \emph{worst} case, with
      $\Theta(1)$ auxiliary space. That combination is unique among the sorts
      in this course.
  \end{itemize}
\end{frame}

\begin{frame}{Heap Sort: The Bill}
  \begin{center}
  {\footnotesize
  \begin{tabular}{lll}
    \toprule
    \textbf{property} & \textbf{heap sort} & \textbf{note} \\
    \midrule
    best / average / worst & $\Theta(n\log n)$ & the guarantee merge sort has \\
    auxiliary space & $\Theta(1)$ & the space quick sort has \\
    stable & no & long-range swaps across the heap \\
    cache behaviour & poor & sift-down jumps by powers of two \\
    \bottomrule
  \end{tabular}}
  \end{center}
  \begin{itemize}
    \item \good{Worst-case guarantee plus bounded memory}: the reason heap
      sort is used in real-time and embedded systems, and as the fallback in
      introsort.
    \item \bad{Not stable}, and \bad{not cache-friendly}: index arithmetic
      jumps across the array, so each sift-down touches memory far from the
      previous access. In practice quick sort is often several times faster
      despite the same $\Theta$ class.
  \end{itemize}
  \muted{The heap itself returns in Week 11 as a priority queue, and in
    Week 12 inside Dijkstra's algorithm --- the structure is worth more than
    this one sort.}
\end{frame}

\begin{frame}{Radix Sort: Stop Comparing Whole Keys}
  \begin{itemize}
    \item The lower bound assumed the only thing a sort may do is compare two
      whole keys. But a $3$-digit number already tells you a lot about where
      it belongs --- the hundreds digit is more significant than the tens.
    \item \key{LSD radix sort} sorts by the least-significant digit first,
      then the next, and so on to the most significant. Each digit pass is a
      \emph{stable} counting sort (Week 8).
  \end{itemize}
  \begin{block}{Definition --- LSD radix sort (Sedgewick Sec.~5.1, p.~702;
    Karumanchi Ch.~10, pp.504--546 PDF)}
    For $d$ digits and a base-$k$ alphabet, run a stable counting sort on
    digit $0$ (units), then digit $1$ (tens), up to digit $d-1$. After pass
    $t$, the array is sorted by its last $t+1$ digits. After the final pass,
    it is fully sorted.
  \end{block}
  \begin{itemize}
    \item The passes must go least-significant first, and each must be
      stable: a later, more-significant pass must not disturb the order the
      earlier passes established within equal digits.
  \end{itemize}
\end{frame}

\begin{frame}{Radix Sort, Drawn: Three Stable Passes}
  \begin{center}
    % Diagram: LSD radix sort of [170,45,75,90,802,24], base 10, 3 digits.
    % Four rows: input, after units, after tens, after hundreds. Cells 1.1 cm
    % wide at x = 0..5; rows y = 2.5, 1.6, 0.7, -0.2.
    \begin{tikzpicture}[every node/.style={font=\scriptsize}]
      \foreach \i/\v in {0/170,1/45,2/75,3/90,4/802,5/24}
        \node[draw, fill=white, minimum width=1.0cm, minimum height=0.55cm,
          inner sep=1pt] at (\i*1.15, 2.5) {\v};
      \node[anchor=east, text=neutral] at (-0.65, 2.5) {input};
      \foreach \i/\v in {0/170,1/90,2/802,3/24,4/45,5/75}
        \node[draw, fill=light-bg, minimum width=1.0cm,
          minimum height=0.55cm, inner sep=1pt] at (\i*1.15, 1.6) {\v};
      \node[anchor=east, text=neutral] at (-0.65, 1.6) {units};
      \foreach \i/\v in {0/802,1/24,2/45,3/75,4/170,5/90}
        \node[draw, fill=light-bg, minimum width=1.0cm,
          minimum height=0.55cm, inner sep=1pt] at (\i*1.15, 0.7) {\v};
      \node[anchor=east, text=neutral] at (-0.65, 0.7) {tens};
      \foreach \i/\v in {0/24,1/45,2/75,3/90,4/170,5/802}
        \node[draw, fill=positive!20, minimum width=1.0cm,
          minimum height=0.55cm, inner sep=1pt] at (\i*1.15, -0.2) {\v};
      \node[anchor=east, text=neutral] at (-0.65, -0.2) {hundreds};
      \node[text=primary-gold, align=center] at (2.9, -1.05)
        {after pass $t$, the last $t+1$ digits are in order};
    \end{tikzpicture}
  \end{center}
  \vspace{-0.1cm}
  \begin{itemize}
    \item Units pass groups by last digit; the tens pass re-groups by tens
      \emph{within} each group; the hundreds pass finishes. Only stability
      keeps the earlier grouping intact.
  \end{itemize}
\end{frame}

\begin{frame}[fragile]{Radix Sort in C: One Counting-Sort Pass per Digit}
  \footnotesize
  \begin{block}{LSD radix sort, base $\le 10$ for clarity (Sedgewick
    Sec.~5.1, p.~702; Karumanchi Ch.~10, pp.504--546 PDF)
    \cite{Sedgewick2011_algorithms,Karumanchi2017_data_structures_algorithms_made_easy}}
\begin{semiverbatim}
void radixSort(int A[], int n, int digits, int base) {
  int *out = malloc(n * sizeof(int));
  for (int d = 0, div = 1; d < digits; d++, div *= base) {
    int count[10] = {0};                 // base <= 10 here
    for (int i = 0; i < n; i++)
      count[(A[i] / div) % base]++;
    for (int v = 1; v < base; v++) count[v] += count[v - 1];
    for (int i = n - 1; i >= 0; i--)     // reverse => stable
      out[--count[(A[i] / div) % base]] = A[i];
    for (int i = 0; i < n; i++) A[i] = out[i];
  }
  free(out);
}
\end{semiverbatim}
  \end{block}
  \vspace{-0.4cm}
  \begin{itemize}
    \item Each pass is Week 8's counting sort on one digit: tally, prefix
      sums, stable reverse scatter. The outer loop advances the digits.
  \end{itemize}
\end{frame}

\begin{frame}{Radix: Beating a Bound by Leaving Its Assumptions}
  \begin{center}
  {\footnotesize
  \begin{tabular}{lll}
    \toprule
    \textbf{cost} & \textbf{value} & \textbf{note} \\
    \midrule
    time & $\Theta(d(n+k))$ & $d$ digits, base $k$ \\
    space & $\Theta(n+k)$ & output buffer + tally \\
    32-bit keys, 8-bit digits & $d=4$, $k=256$ & four linear passes \\
    \bottomrule
  \end{tabular}}
  \end{center}
  \begin{itemize}
    \item It makes \emph{zero} comparisons of whole keys, so the
      $\Omega(n\log n)$ comparison bound simply does not apply. For fixed
      width keys with a small base, $\Theta(d(n+k))$ is linear in $n$.
    \item \good{Stable}; \bad{not in-place} ($\Theta(n+k)$ auxiliary).
    \item The precondition is real: keys must have a positional
      representation with a known width --- integers, fixed-length strings.
      Arbitrary objects with only a comparison function cannot use it.
  \end{itemize}
  \begin{exampleblock}{The lesson}
    A lower bound is only as strong as its model. Radix sort does not break
    the mathematics --- it changes the question from ``compare two keys'' to
    ``read one digit''.
  \end{exampleblock}
\end{frame}

\begin{frame}{Shell Sort: Insertion Sort with a Head Start}
  \begin{itemize}
    \item Insertion sort is $O(n)$ on nearly-sorted input and $\Theta(n^2)$
      on random input. Shell sort manufactures ``nearly sorted'' cheaply: it
      insertion-sorts widely spaced subsequences first, then narrows the
      spacing.
  \end{itemize}
  \begin{block}{Definition --- diminishing-increment sort
    (Wirth Sec.~2.3.1, p.~52; Sedgewick Sec.~2.1, p.~244)}
    Choose a decreasing \key{gap sequence} $h_1 > h_2 > \cdots > h_t = 1$.
    For each gap $h$, insertion-sort every subsequence
    $A[i], A[i+h], A[i+2h], \dots$. The final pass with $h=1$ is ordinary
    insertion sort, so correctness is automatic.
  \end{block}
  \begin{itemize}
    \item Why it helps: a gap of $n/2$ can move an element halfway across the
      array in one shift. By the time $h=1$, few inversions remain, so the
      final insertion pass is close to linear.
    \item Shell's original sequence $n/2, n/4, \dots, 1$ is simple but has a
      $\Theta(n^2)$ worst case. Knuth's $(3^k-1)/2$ sequence gives
      $O(n^{3/2})$; Sedgewick's gives $O(n^{4/3})$. No sequence is known to
      be optimal.
  \end{itemize}
\end{frame}

\begin{frame}{Shell Sort, Drawn: Gaps $4$, then $2$, then $1$}
  \begin{center}
    % Diagram: shell sort of [7,3,9,1,8,2,6,4] with gaps 4, 2, 1. Four rows
    % at y = 2.4, 1.5, 0.6, -0.3; cells 0.85 cm wide at x = 0..7. Colours
    % group the subsequences touched by each gap pass.
    \begin{tikzpicture}[every node/.style={font=\scriptsize}]
      \foreach \i/\v in {0/7,1/3,2/9,3/1,4/8,5/2,6/6,7/4}
        \node[draw, fill=white, minimum width=0.75cm, minimum height=0.5cm,
          inner sep=0] at (\i, 2.4) {\v};
      \node[anchor=east, text=neutral] at (-0.55, 2.4) {input};
      % gap 4 result
      \foreach \i/\v/\c in {0/7/1,1/2/2,2/6/3,3/1/4,4/8/1,5/3/2,6/9/3,7/4/4}
        \node[draw, fill=primary-gold!\c0, minimum width=0.75cm,
          minimum height=0.5cm, inner sep=0] at (\i, 1.5) {\v};
      \node[anchor=east, text=neutral] at (-0.55, 1.5) {gap $4$};
      % gap 2 result: two interleaved groups
      \foreach \i/\v in {0/6,1/1,2/7,3/2,4/8,5/3,6/9,7/4}
        \node[draw, fill=light-bg, minimum width=0.75cm,
          minimum height=0.5cm, inner sep=0] at (\i, 0.6) {\v};
      \foreach \i in {0,2,4,6}
        \node[text=positive] at (\i, 0.95) {\tiny $\bullet$};
      \foreach \i in {1,3,5,7}
        \node[text=negative] at (\i, 0.95) {\tiny $\bullet$};
      \node[anchor=east, text=neutral] at (-0.55, 0.6) {gap $2$};
      % gap 1 result
      \foreach \i/\v in {0/1,1/2,2/3,3/4,4/6,5/7,6/8,7/9}
        \node[draw, fill=positive!15, minimum width=0.75cm,
          minimum height=0.5cm, inner sep=0] at (\i, -0.3) {\v};
      \node[anchor=east, text=neutral] at (-0.55, -0.3) {gap $1$};
      \node[text=primary-gold, align=center] at (3.5, -1.05)
        {each pass is insertion sort on one residue class mod the gap};
    \end{tikzpicture}
  \end{center}
  \vspace{-0.1cm}
  \begin{itemize}
    \item Row 2 groups positions $\{0,4\}$, $\{1,5\}$, $\{2,6\}$,
      $\{3,7\}$; row 3 groups the even and odd positions separately; row 4 is
      plain insertion sort.
  \end{itemize}
\end{frame}

\begin{frame}[fragile]{Shell Sort in C}
  \scriptsize
  \begin{block}{Shell sort with the halving sequence (Wirth Sec.~2.3.1,
    p.~52; Karumanchi Ch.~10, pp.504--546 PDF)
    \cite{Wirth2004_algorithms_data_structures,Karumanchi2017_data_structures_algorithms_made_easy}}
\begin{semiverbatim}
void shellSort(int A[], int n) {
  for (int gap = n / 2; gap > 0; gap /= 2) {  // Shell's halving sequence
    for (int i = gap; i < n; i++) {
      int key = A[i], j = i;
      while (j >= gap && A[j - gap] > key) {  // shifted insertion
        A[j] = A[j - gap];
        j -= gap;
      }
      A[j] = key;
    }
  }
}
\end{semiverbatim}
  \end{block}
  \vspace{-0.3cm}
  \begin{itemize}
    \item Compare with Week 8's insertion sort: the only change is the stride
      \texttt{gap} in place of $1$; with \texttt{gap = 1} the code is exactly
      the old insertion sort, which is the correctness argument.
    \item \good{In-place}, \bad{unstable}, simpler than quick sort --- but
      the worst case depends on the gap sequence, so it carries no guarantee.
  \end{itemize}
\end{frame}

\begin{frame}{Check Your Prediction: Choose the Sort}
  \begin{block}{Match each workload to the best sort, and name the property
    that decides it.}
    (i) $10^7$ 32-bit integers, memory tight, worst case must be guaranteed.
    \quad (ii) $10^6$ 16-bit IDs, base-256 digits available.
    \quad (iii) $2{,}000$ records already almost sorted.
    \quad (iv) a general-purpose library sort facing adversarial input.
  \end{block}
  \begin{enumerate}
    \item (i) \key{heap sort}: $\Theta(n\log n)$ worst case with $O(1)$
      auxiliary space --- the only sort with both.
    \item (ii) \key{radix sort}: two base-256 passes, $\Theta(2(n+256))$,
      linear in $n$; the keys are fixed-width integers.
    \item (iii) \key{insertion sort} (or shell sort): nearly sorted input is
      insertion sort's $O(n)$ best case.
    \item (iv) \key{quick sort} with a randomized pivot, plus a heap-sort
      fallback (introsort): fast average, and the guarantee restored when
      recursion gets too deep.
  \end{enumerate}
\end{frame}

\transitionslide{Five sorts, one wall, and a decision to make.}

% =============================================================================
% ACT 5 --- COMPARATIVE ANALYSIS AND THE CLOSE
% =============================================================================
\section{Comparative Analysis and Close}

\sectiondivider{Section 5}{Comparative Analysis and the Close}

\begin{frame}{The Efficient Sorts, Side by Side}
  \begin{center}
  {\scriptsize
  \begin{tabular}{l l l l l c c}
    \toprule
    \textbf{sort} & \textbf{best} & \textbf{average} & \textbf{worst} &
      \textbf{space} & \textbf{stable} & \textbf{in-place} \\
    \midrule
    merge & $\Theta(n\log n)$ & $\Theta(n\log n)$ & $\Theta(n\log n)$ &
      $\Theta(n)$ & yes & no \\
    quick & $\Theta(n\log n)$ & $\Theta(n\log n)$ & $\Theta(n^2)$ &
      $O(\log n)$ & no & yes \\
    heap & $\Theta(n\log n)$ & $\Theta(n\log n)$ & $\Theta(n\log n)$ &
      $O(1)$ & no & yes \\
    shell & $O(n\log n)$ & depends on gaps & $\Theta(n^2)$ / $O(n^{3/2})$ &
      $O(1)$ & no & yes \\
    radix & $\Theta(d(n+k))$ & $\Theta(d(n+k))$ & $\Theta(d(n+k))$ &
      $\Theta(n+k)$ & yes & no \\
    \bottomrule
  \end{tabular}}
  \end{center}
  \vspace{0.1cm}
  \begin{itemize}
    \item The shell row's two worst-case entries are Shell's halving sequence
      ($\Theta(n^2)$) and Knuth's $(3^k-1)/2$ sequence ($O(n^{3/2})$); the
      cost is a property of the gap sequence, not of ``shell sort'' alone.
    \item Read the table by \emph{which cell is guaranteed}. Only merge and
      heap guarantee $\Theta(n\log n)$; only heap combines that with
      $O(1)$ space.
  \end{itemize}
\end{frame}

\begin{frame}{Choosing: Read the Workload, Then the Sort}
  \begin{center}
  {\footnotesize
  \begin{tabular}{p{6.6cm} p{5.2cm}}
    \toprule
    \textbf{workload} & \textbf{choose} \\
    \midrule
    worst-case guarantee + tight memory & heap sort \\
    raw average speed, in-place, cache-friendly & quick sort (randomized) \\
    stability needed, or linked / external data & merge sort \\
    fixed-width integer or string keys, known range & radix sort \\
    nearly sorted, or small $n$ & insertion sort \\
    simple in-place, no guarantee required & shell sort \\
    \bottomrule
  \end{tabular}}
  \end{center}
  \begin{itemize}
    \item No single winner. Every row is correct for \emph{some} input, and
      the exam question is always ``state the workload, then justify the
      choice'' --- never ``rank the sorts''.
    \item The three properties that decide it are the same three Week 8
      introduced: the case, stability, and space. Week 9 only adds the
      guarantee that a $\Theta(n\log n)$ worst case is even possible.
  \end{itemize}
\end{frame}

\begin{frame}{Three Classic Mistakes}
  \begin{itemize}
    \item \bad{``Quick sort is $O(n\log n)$'' with no case.} Its worst case
      is $\Theta(n^2)$ --- and the input that triggers it (sorted data with a
      last-element pivot) is exactly the input a tokenizer produces. State
      the case, then name the pivot rule.
    \item \bad{``Merge sort is in-place'' / ``quick sort is $O(1)$ space''.}
      Merge sort needs a $\Theta(n)$ buffer; quick sort needs $O(\log n)$
      recursion stack (unless tail-call optimisation removes it). Heap sort
      is the one efficient sort that is genuinely $O(1)$ auxiliary.
    \item \bad{``Radix sort is $O(n)$, so it always beats merge sort.''} It is
      $\Theta(d(n+k))$, linear only when $d$ and $k$ are treated as
      constants, and it needs $\Theta(n+k)$ space and fixed-width keys.
      Neither statement is true for arbitrary comparable objects.
  \end{itemize}
  \begin{exampleblock}{How to revise}
    For each of the five sorts, write one line: the idea, the recurrence or
    pass count, the worst case, stability, and space. Then trace a partition
    and a heap extraction by hand. The test is traces and cases, not names.
  \end{exampleblock}
\end{frame}

\begin{frame}{Week 9 Summary}
  \begin{itemize}
    \item \key{Divide and conquer} --- merge sort splits, sorts, and merges.
      $T(n)=2T(n/2)+\Theta(n)=\Theta(n\log n)$ in every case; stable; needs
      $\Theta(n)$ space.
    \item \key{Partition} --- quick sort places a pivot and recurses.
      $\Theta(n\log n)$ average, $\Theta(n^2)$ worst on sorted input with a
      last-element pivot; in-place and unstable; randomize or use
      median-of-three.
    \item \key{The wall} --- any comparison sort needs $\Omega(n\log n)$
      comparisons, because a decision tree must distinguish all $n!$ orders.
      Merge and heap reach the bound; they cannot beat it.
    \item \key{Three ways past it} --- heap sort keeps the guarantee in
      $O(1)$ space; radix sort leaves the comparison model with stable digit
      passes, $\Theta(d(n+k))$; shell sort buys simplicity and in-place
      operation with a gap-dependent bound.
  \end{itemize}
  \begin{exampleblock}{Next week}
    The table is now sorted and searchable in $O(\log n)$. Week 10 asks a
    stronger question: can we find a record in \emph{one} probe, without
    sorting at all? File organization and hashing.
  \end{exampleblock}
  \vspace{-0.2cm}
\end{frame}

\begin{frame}{Before You Go: Predict}
  \begin{block}{Three one-line predictions}
    Answer from the algorithms, then check in the lab.
  \end{block}
  \begin{enumerate}
    \item Merge sort on $8$ elements: how many levels of merging, and does an
      already-sorted input reduce the number of comparisons?
    \item Quick sort $[5,4,3,2,1]$ with the last-element pivot: how many
      comparisons does \texttt{partition} perform in total?
    \item You must sort $10^6$ 16-bit student IDs with memory to spare. Which
      sort, and how many digit passes at base $256$?
  \end{enumerate}
  \muted{Answers: $\log_2 8 = 3$ merging levels; no --- sorted input still
    costs $\Theta(n\log n)$ comparisons, because every merge walk is
    unchanged. $4+3+2+1 = 10$ comparisons --- the reverse-sorted worst case.
    Radix sort: $16$ bits $/\,8$ bits per digit $= 2$ passes, each linear,
    with $\Theta(n+256)$ extra space.}
\end{frame}

\begin{frame}[allowframebreaks]{References}
  \tiny
  \bibliographystyle{plain}
  \bibliography{../../Bibliography_base}
\end{frame}

\end{document}
